\documentclass[t, aspectratio=169]{beamer}
\usepackage[utf8]{vietnam}
\usepackage{amsthm,amsmath,amssymb}
\usepackage{lmodern}
\usepackage{listings}
\usepackage{xcolor}

% --- CẤU HÌNH LỀ ---
\def\slideContentLeftMargin{0.5cm}
\def\slideContentRightMargin{1cm}
\def\hustHeaderHeight{1cm}
\def\slideContentBotMargin{3.5cm}

% Sử dụng theme HUST màu xanh, tỉ lệ 16:9
\usepackage[blue, 169]{beamerthemeHUST} 

% Điều chỉnh vị trí title page
\renewcommand{\titlePageTitleX}{-0.5cm}
\renewcommand{\titlePageTitleY}{-0.5cm}

% Căn lề cho trang title
\renewcommand{\hustTitleAlign}{\alignLeft} 
\renewcommand{\hustSubtitleAlign}{\alignLeft}
\renewcommand{\hustAuthorAlign}{\alignLeft}
\renewcommand{\hustInstAlign}{\alignLeft}

% Cấu hình listings cho code C
\lstset{
    language=C,
    basicstyle=\tiny\ttfamily,
    keywordstyle=\color{blue}\bfseries,
    commentstyle=\color{green!60!black},
    stringstyle=\color{red},
    numbers=left,
    numberstyle=\tiny\color{gray},
    stepnumber=1,
    numbersep=5pt,
    backgroundcolor=\color{gray!10},
    showspaces=false,
    showstringspaces=false,
    showtabs=false,
    frame=single,
    tabsize=2,
    captionpos=b,
    breaklines=true,
    breakatwhitespace=false,
    escapeinside={(*@}{@*)},
    xleftmargin=0.5cm,
    xrightmargin=0.5cm
}

% Thông tin slide
\title{GIẢI THUẬT VÀ CÀI ĐẶT CHƯƠNG TRÌNH SONG SONG \\ BÀI TOÁN CÓ SỰ PHỤ THUỘC DỮ LIỆU}
\subtitle{Lattice Boltzmann Method (D2Q9)}
\author{Nhóm 1}
\institute{IT5408 - Tính toán Hiệu năng Cao}
%\date{2026}

\begin{document}
	
	\hustintropage{}
	\hustsectionpage{}
	
	
	
	% Trang bìa
	\husttitlepage
	
	% Slide thành viên Nhóm 1
	\begin{frame}{Thành viên Thực hiện - Nhóm 1}
		\begin{table}
			\centering
			\scriptsize
			\begin{tabular}{|c|l|c|p{6.8cm}|}
				\hline
				\textbf{STT} & \textbf{Họ và tên} & \textbf{Mã HV} & \textbf{Phân công nhiệm vụ chính} \\
				\hline
				1 & \textbf{Phạm Tiến Thành} & \textbf{20242422M} & Trưởng nhóm, Cài đặt song song MPI, Benchmark \\
				\hline
				2 & Trần Thị Hạnh & 20242058M & Phân tích mô hình toán D2Q9, Điều kiện biên vật lý \\
				\hline
				3 & Bùi Minh Đức & 20242056M & Cài đặt thuật toán tuần tự C99, Kiểm thử Checksum \\
				\hline
				4 & Nguyễn Thị Huê & 20252584M & Xử lý dữ liệu mô phỏng, Trực quan hóa kết quả \\
				\hline
				5 & Nguyễn Đình Hiếu & 20261223M & Thiết kế Slide Beamer, Phân tích Scaling hiệu năng \\
				\hline
				6 & Đỗ Minh Hiếu & 20242031M & Biên soạn báo cáo kỹ thuật, Phân tích Overhead \\
				\hline
			\end{tabular}
		\end{table}
	\end{frame}

	
	% Slide thành viên Nhóm 1
	\begin{frame}{Thành viên Thực hiện - Nhóm 1}
		\begin{table}
			\centering
			\scriptsize
			\begin{tabular}{|c|l|c|p{6.8cm}|}
				\hline
				\textbf{STT} & \textbf{Họ và tên} & \textbf{Mã HV} & \textbf{Phân công nhiệm vụ chính} \\
				\hline
				1 & \textbf{Phạm Tiến Thành} & \textbf{20242422M} & Trưởng nhóm, Cài đặt song song MPI, Benchmark \\
				\hline
				2 & Trần Thị Hạnh & 20242058M & Phân tích mô hình toán D2Q9, Điều kiện biên vật lý \\
				\hline
				3 & Bùi Minh Đức & 20242056M & Cài đặt thuật toán tuần tự C99, Kiểm thử Checksum \\
				\hline
				4 & Nguyễn Thị Huê & 20252584M & Xử lý dữ liệu mô phỏng, Trực quan hóa kết quả \\
				\hline
				5 & Nguyễn Đình Hiếu & 20261223M & Thiết kế Slide Beamer, Phân tích Scaling hiệu năng \\
				\hline
				6 & Đỗ Minh Hiếu & 20242031M & Biên soạn báo cáo kỹ thuật, Phân tích Overhead \\
				\hline
			\end{tabular}
		\end{table}
	\end{frame}

	
	\begin{frame}[allowframebreaks]{Nội dung trình bày}
		\tableofcontents
	\end{frame}
	
	\resetPageCount
	\showPageNumber
	
	% =================================================================
	% PHẦN 1: MÔ HÌNH TOÁN HỌC
	% =================================================================
	\useStyleHustFull
	\useThinBar
	
	\section{Mô hình Toán học và Phương pháp Giải}
	
	\begin{frame}{Phương trình Lattice Boltzmann}
		\textbf{Phương trình LBE (Lattice Boltzmann Equation):}
		\begin{equation}
			f_i(x + c_i \Delta t, t + \Delta t) = f_i(x, t) - \omega[f_i(x, t) - f_i^{eq}(x, t)]
		\end{equation}
		
		\vspace{0.5em}
		
		\textbf{Công thức giải:}
		\begin{itemize}
			\item \textbf{Khởi tạo:} $f_i(x, 0) = f_i^{eq}(\rho=1, u=0)$
			\item \textbf{Tại bước thời gian $n+1$:}
			\begin{enumerate}
				\item \textbf{Collision (Va chạm):} 
				\[f^*_i(x, t) = f_i(x, t) - \omega(f_i(x, t) - f_i^{eq}(x, t))\]
				\item \textbf{Streaming (Di chuyển):} 
				\[f_i(x + c_i \Delta t, t + \Delta t) = f^*_i(x, t)\]
			\end{enumerate}
		\end{itemize}
	\end{frame}
	
	\begin{frame}{Mô hình D2Q9 Stencil}
		\begin{columns}[T]
			\begin{column}{0.5\textwidth}
				\textbf{Mô hình D2Q9:}
				\begin{itemize}
					\item \textbf{D2:} 2 chiều (2D)
					\item \textbf{Q9:} 9 hướng vận tốc
					\item Vector vận tốc: $c_i = (c_{ix}, c_{iy})$
					\item Trọng số: $w_i$
				\end{itemize}
				
				\vspace{1em}
				
				\textbf{Sơ đồ hướng:}
				\begin{center}
					\begin{tabular}{ccc}
						6 & 2 & 5 \\
						$\nwarrow$ & $\uparrow$ & $\nearrow$ \\
						3 & 0 & 1 \\
						$\leftarrow$ & $\bullet$ & $\rightarrow$ \\
						7 & 4 & 8 \\
						$\swarrow$ & $\downarrow$ & $\searrow$ \\
					\end{tabular}
				\end{center}
			\end{column}
			
			\begin{column}{0.5\textwidth}
				\textbf{Trọng số và Vector:}
				\begin{itemize}
					\item $w_0 = \frac{4}{9}$ (center)
					\item $w_{1-4} = \frac{1}{9}$ (cardinal)
					\item $w_{5-8} = \frac{1}{36}$ (diagonal)
				\end{itemize}
				
				\vspace{1em}
				
				\textbf{Vector vận tốc:}
				\begin{itemize}
					\item $c_0 = (0, 0)$
					\item $c_{1,2,3,4} = (\pm1, 0), (0, \pm1)$
					\item $c_{5,6,7,8} = (\pm1, \pm1)$
				\end{itemize}
			\end{column}
		\end{columns}
	\end{frame}
	
	% =================================================================
	% PHẦN 2: SỰ PHỤ THUỘC DỮ LIỆU
	% =================================================================
	\section{Sự Phụ Thuộc Dữ Liệu}
	
	\begin{frame}{Phân tích Phụ thuộc Dữ liệu}
		\textbf{Sự phụ thuộc xảy ra ở bước Streaming:}
		\begin{itemize}
			\item Tại vị trí $x$, hàm phân bố $f_i$ ở thời điểm tiếp theo phụ thuộc vào giá trị sau va chạm $f^*_i$ từ điểm láng giềng.
			\item Cụ thể: $f_i(x, t+1)$ nhận giá trị từ điểm $x - c_i \Delta t$
		\end{itemize}
		
		\vspace{1em}
		
		\textbf{Xử lý trong chương trình:}
		\begin{block}{Tuần tự}
			\begin{itemize}
				\item Điều kiện biên Bounce-back tại tường (Y)
				\item Điều kiện biên Periodic tại biên X
			\end{itemize}
		\end{block}
		
		\begin{alertblock}{Song song}
			\begin{itemize}
				\item Cần trao đổi hàm phân bố biên giữa các CPU
				\item Sử dụng \textbf{Ghost Columns} (Cột ảo)
				\item Trao đổi qua MPI trước Streaming
			\end{itemize}
		\end{alertblock}
	\end{frame}
	
	% =================================================================
	% PHẦN 3: ĐẠI LƯỢNG VĨ MÔ
	% =================================================================
	\section{Đại Lượng Vĩ Mô và Ký Hiệu}
	
	\begin{frame}{Các Đại lượng Vĩ mô}
		\textbf{Ký hiệu:}
		\begin{itemize}
			\item $f_i$: Hàm phân bố hạt theo hướng $i$ ($i=0 \dots 8$)
			\item $c_i$: Vector vận tốc rời rạc hướng $i$
			\item $w_i$: Trọng số lý thuyết
			\item $\omega = 1/\tau$: Tham số thư giãn
		\end{itemize}
		
		\vspace{1em}
		
		\begin{columns}[T]
			\begin{column}{0.5\textwidth}
				\textbf{Đại lượng vĩ mô:}
				\begin{itemize}
					\item \textbf{Mật độ:} 
					\[\rho = \sum_{i=0}^{8} f_i\]
					\item \textbf{Vận tốc:} 
					\[u = \frac{1}{\rho} \sum_{i=0}^{8} f_i c_i\]
				\end{itemize}
			\end{column}
			
			\begin{column}{0.5\textwidth}
				\textbf{Hàm phân bố cân bằng:}
				\begin{equation}
					\begin{split}
						f_i^{eq} = w_i \rho \Big(1 + 3(c_i \cdot u) \\
						+ \frac{9}{2}(c_i \cdot u)^2 - \frac{3}{2}u^2\Big)
					\end{split}
				\end{equation}
			\end{column}
		\end{columns}
	\end{frame}
	
	% =================================================================
	% PHẦN 4: CÀI ĐẶT CODE
	% =================================================================
	\useStyleLogoOnly
	\useThickBar
	
	\section{Cài Đặt Thực Tế}
	
	\begin{frame}[fragile]{Code: Collision và TinhMacro}
		\begin{lstlisting}[caption={Hàm Collision - Toán tử BGK}]
void CollisionCucBo(double *f, double *f_new, double *rho, 
                     double *ux, double *uy, int nx_local, 
                     int ny, double omega) {
    for (int x = 0; x < nx_local; x++) {
        for (int y = 0; y < ny; y++) {
            double r = *(rho + x*ny + y);
            double u_x = *(ux + x*ny + y);
            double u_y = *(uy + x*ny + y);
            double usqr = u_x*u_x + u_y*u_y;
            
            for (int i = 0; i < Q; i++) {
                double cu = cx[i]*u_x + cy[i]*u_y;
                double feq = w[i] * r * (1.0 + 3.0*cu 
                                    + 4.5*cu*cu - 1.5*usqr);
                double fi = *(f + (x*ny + y)*Q + i);
                (*@\colorbox{yellow}{$*(f\_new + (x*ny + y)*Q + i) = fi - omega * (fi - feq);$}@*)
            }
        }
    }
}
		\end{lstlisting}
	\end{frame}
	
	\begin{frame}[fragile]{Code: Streaming với Ghost Columns}
		\begin{lstlisting}[caption={Hàm StreamingCucBo - Xử lý biên với Ghost}]
void StreamingCucBo(double *f_new, double *f, int nx_local, 
                     int ny, double *ghost_left, 
                     double *ghost_right, int rank, int size) {
    for (int x = 0; x < nx_local; x++) {
        for (int y = 0; y < ny; y++) {
            for (int i = 0; i < Q; i++) {
                int xn = x + cx[i]; 
                int yn = y + cy[i];
                
                if (yn < 0 || yn >= ny) { // Bounce-back
                    *(f + (x*ny+y)*Q + opposite[i]) += *(f_new + (x*ny+y)*Q + i);
                } else if (xn < 0) { // Ghost Left
                    (*@\colorbox{yellow}{if (rank > 0) *(f + (x*ny+yn)*Q + i) += *(ghost\_left + yn*Q + i);}@*)
                } else if (xn >= nx_local) { // Ghost Right
                    (*@\colorbox{yellow}{if (rank < size-1) *(f + (x*ny+yn)*Q + i) += *(ghost\_right + yn*Q + i);}@*)
                } else { // Trong miền
                    *(f + (xn*ny+yn)*Q + i) += *(f_new + (x*ny+y)*Q + i);
                }
            }
        }
    }
}
		\end{lstlisting}
	\end{frame}
	
	% =================================================================
	% PHẦN 5: GIẢI THUẬT SONG SONG
	% =================================================================
	\section{Giải Thuật Song Song SPMD}
	
	\begin{frame}{Chiến lược Song song hóa}
		\textbf{Mô hình:} Single Program Multiple Data (SPMD)
		
		\vspace{1em}
		
		\begin{columns}[T]
			\begin{column}{0.5\textwidth}
				\textbf{Chia miền (Domain Decomposition):}
				\begin{itemize}
					\item Chia 1D theo cột (chiều X)
					\item Mỗi tiến trình: $nx\_local \approx NX/P$ cột
					\item Chiều Y: Không chia
				\end{itemize}
				
				\vspace{0.5em}
				
				\textbf{Ví dụ: 256×64 với 4 tiến trình}
				\begin{itemize}
					\item Rank 0: Cột 0-63
					\item Rank 1: Cột 64-127
					\item Rank 2: Cột 128-191
					\item Rank 3: Cột 192-255
				\end{itemize}
			\end{column}
			
			\begin{column}{0.5\textwidth}
				\textbf{Vùng đệm biên:}
				\begin{itemize}
					\item \alert{Ghost Columns}
					\item Kích thước: $NY \times Q$ doubles
					\item Với $NY=64$: $64 \times 9 \times 8 = 4.5$ KB
				\end{itemize}
				
				\vspace{0.5em}
				
				\textbf{Giao tiếp:}
				\begin{itemize}
					\item Trao đổi với láng giềng trái/phải
					\item Sử dụng MPI non-blocking
					\item \texttt{MPI\_Isend} / \texttt{MPI\_Irecv}
					\item \texttt{MPI\_Waitall} đồng bộ
				\end{itemize}
			\end{column}
		\end{columns}
	\end{frame}
	
	\begin{frame}{5 Bước Chính trong Giải thuật}
		\begin{enumerate}
			\item \textbf{Khởi tạo:}
			\begin{itemize}
				\item Cấp phát bộ nhớ: $f$, $f_{new}$, $\rho$, $u$, ghost buffers
				\item Đặt trạng thái ban đầu: $\rho=1$, $u=0$
			\end{itemize}
			
			\item \textbf{Chia miền:}
			\begin{itemize}
				\item Tính $nx\_local$ cho mỗi Rank
				\item Xử lý phần dư: $NX \bmod P$
			\end{itemize}
			
			\item \textbf{Phân tán:}
			\begin{itemize}
				\item Mỗi Rank tự khởi tạo vùng của mình
				\item Hoặc Rank 0 phát tán dữ liệu
			\end{itemize}
			
			\item \textbf{Lặp tính toán:} (10,000 bước)
			\begin{itemize}
				\item Tính đại lượng vĩ mô $\to$ Collision $\to$ \alert{Trao đổi biên} $\to$ Streaming
			\end{itemize}
			
			\item \textbf{Tổng hợp:}
			\begin{itemize}
				\item Rank 0 thu thập kết quả (\texttt{MPI\_Reduce})
				\item Ghi file output
			\end{itemize}
		\end{enumerate}
	\end{frame}
	
	% =================================================================
	% PHẦN 6: TRUYỀN THÔNG MPI
	% =================================================================
	\section{Chi Tiết Truyền Thông}
	
	\begin{frame}{Trao đổi Ghost Columns}
		\textbf{Sử dụng Non-blocking Communication:}
		\begin{itemize}
			\item Tối ưu hóa thời gian chờ
			\item Cho phép tính toán và giao tiếp đồng thời
		\end{itemize}
		
		\vspace{1em}
		
		\textbf{Quy trình trao đổi:}
		\begin{enumerate}
			\item Rank $i$ gửi cột 0 sang Rank $i-1$, nhận vào \texttt{ghost\_left}
			\item Rank $i$ gửi cột cuối sang Rank $i+1$, nhận vào \texttt{ghost\_right}
			\item Sử dụng \texttt{MPI\_Waitall} để đồng bộ
		\end{enumerate}
		
		\vspace{1em}
		
		\begin{alertblock}{Quan trọng}
			Phải đợi tất cả giao tiếp hoàn thành trước khi thực hiện Streaming!
		\end{alertblock}
	\end{frame}
	
	\begin{frame}[fragile]{Code MPI: TraoDoiGhost}
		\begin{lstlisting}[caption={Hàm trao đổi Ghost Columns}]
void TraoDoiGhost(double *f_new, int nx_local, int ny, 
                   double *ghost_left, double *ghost_right, 
                   int rank, int size, MPI_Comm comm) {
    MPI_Request reqs[4]; 
    int rcount = 0;
    
    (*@\colorbox{yellow}{// Trao đổi với Rank trái}@*)
    if (rank > 0) {
        MPI_Irecv(ghost_left, ny*Q, MPI_DOUBLE, 
                  rank-1, 100, comm, &reqs[rcount++]);
        MPI_Isend(f_new, ny*Q, MPI_DOUBLE, 
                  rank-1, 101, comm, &reqs[rcount++]);
    }
    
    (*@\colorbox{yellow}{// Trao đổi với Rank phải}@*)
    if (rank < size-1) {
        MPI_Irecv(ghost_right, ny*Q, MPI_DOUBLE, 
                  rank+1, 101, comm, &reqs[rcount++]);
        MPI_Isend(f_new + (nx_local-1)*ny*Q, ny*Q, MPI_DOUBLE, 
                  rank+1, 100, comm, &reqs[rcount++]);
    }
    
    if (rcount > 0) MPI_Waitall(rcount, reqs, MPI_STATUSES_IGNORE);
}
		\end{lstlisting}
	\end{frame}
	
	% =================================================================
	% PHẦN 7: KẾT QUẢ THỰC NGHIỆM
	% =================================================================
	\useThickBar
	\useStyleLogoOnly
	
	\section{Kết Quả Thực Nghiệm}
	
	\begin{frame}{Hiệu năng - Lưới 256×64}
		\textbf{Thông số:} $\omega=1.0$, $Steps=10{,}000$, $\nu = 0.1667$
		
		\vspace{1em}
		
		\begin{columns}[T]
			\begin{column}{0.55\textwidth}
				\begin{block}{Kết quả đo đạc}
					\centering
					\begin{tabular}{|c|c|c|c|c|}
						\hline
						\textbf{CPU} & \textbf{Time (s)} & \textbf{MLUPS} & \textbf{Speedup} & \textbf{Eff.} \\
						\hline
						1 & 1.802 & 90.94 & 1.00× & 100\% \\
						\hline
						2 & 1.365 & 120.00 & 1.32× & 66\% \\
						\hline
						\alert{4} & \alert{0.738} & \alert{222.13} & \alert{2.44×} & \alert{61\%} \\
						\hline
						8 & 0.450 & 364.09 & 4.00× & 50\% \\
						\hline
					\end{tabular}
				\end{block}
				
				\vspace{0.5em}
				\footnotesize
				\textit{*MLUPS: Million Lattice Updates Per Second}
			\end{column}
			
			\begin{column}{0.45\textwidth}
				\textbf{Nhận xét:}
				\begin{itemize}
					\item Speedup tốt với 2-4 tiến trình
					\item Hiệu suất $>60\%$ với 4 CPUs
					\item MLUPS tăng gần tuyến tính
					\item Overhead tăng với nhiều tiến trình
				\end{itemize}
				
				\vspace{1em}
				
				\textbf{Nguyên nhân giảm hiệu suất:}
				\begin{itemize}
					\item Giao tiếp MPI
					\item Đồng bộ hóa
					\item Cache miss
				\end{itemize}
			\end{column}
		\end{columns}
	\end{frame}
	
	\begin{frame}{Độ chính xác và Visualization}
		\begin{columns}[T]
			\begin{column}{0.5\textwidth}
				\textbf{Xác thực độ chính xác:}
				\begin{itemize}
					\item Vận tốc trung bình: Serial vs MPI
					\item Sai số: $< 10^{-14}$ (machine epsilon)
					\item \alert{Kết quả khớp hoàn toàn!}
				\end{itemize}
				
				\vspace{1em}
				
				\textbf{Trường vận tốc:}
				\begin{itemize}
					\item Đạt trạng thái ổn định
					\item Profile parabol chuẩn
					\item Phù hợp lý thuyết Poiseuille
				\end{itemize}
				
				\vspace{1em}
				
				\textbf{Bảo toàn khối lượng:}
				\begin{itemize}
					\item $\rho_{total} = const$
					\item Sai số $< 0.01\%$
				\end{itemize}
			\end{column}
			
			\begin{column}{0.5\textwidth}
				\begin{figure}
					\centering
					\fbox{
						\begin{minipage}[c][4cm]{0.9\textwidth}
							\centering 
							\vspace{1.5cm}
							\textbf{[Velocity Field Plot]}
							\vspace{1cm}
							
							\footnotesize
							\textit{Trường vận tốc với profile Poiseuille}
						\end{minipage}
					}
				\end{figure}
				
				\begin{figure}
					\centering
					\fbox{
						\begin{minipage}[c][2.5cm]{0.9\textwidth}
							\centering 
							\vspace{0.5cm}
							\textbf{[Performance Charts]}
							\vspace{0.5cm}
							
							\footnotesize
							\textit{Biểu đồ Speedup và Efficiency}
						\end{minipage}
					}
				\end{figure}
			\end{column}
		\end{columns}
	\end{frame}
	
	% =================================================================
	% PHẦN 8: KẾT LUẬN
	% =================================================================
	\section{Kết Luận}
	
	\begin{frame}{Kết luận và Hướng phát triển}
		\begin{columns}[T]
			\begin{column}{0.5\textwidth}
				\textbf{Ưu điểm của LBM:}
				\begin{itemize}
					\item Tính cục bộ cao
					\item Dễ song song hóa
					\item Xử lý biên phức tạp đơn giản
					\item Phù hợp với MPI/GPU
				\end{itemize}
				
				\vspace{1em}
				
				\textbf{Hạn chế:}
				\begin{itemize}
					\item Overhead truyền thông
					\item Hiệu suất giảm với lưới nhỏ
					\item Yêu cầu bộ nhớ lớn (9 hướng)
				\end{itemize}
			\end{column}
			
			\begin{column}{0.5\textwidth}
				\textbf{Hướng phát triển:}
				\begin{enumerate}
					\item \textbf{Chia miền 2D:}
					\begin{itemize}
						\item Giảm overhead giao tiếp
						\item Tốt hơn với nhiều CPU
					\end{itemize}
					
					\item \textbf{Overlapping:}
					\begin{itemize}
						\item Tính toán trong khi giao tiếp
						\item Tăng hiệu suất 10-20\%
					\end{itemize}
					
					\item \textbf{GPU (CUDA):}
					\begin{itemize}
						\item Speedup 50-100×
						\item Phù hợp cho lưới lớn
					\end{itemize}
					
					\item \textbf{Hybrid MPI+OpenMP:}
					\begin{itemize}
						\item Tận dụng đa nhân
						\item Giảm số message MPI
					\end{itemize}
				\end{enumerate}
			\end{column}
		\end{columns}
	\end{frame}
	
	% =================================================================
	% THAM KHẢO
	% =================================================================
	\useStyleLogoOnly
	
	\setbeamertemplate{bibliography item}[text]
	
	\begin{frame}[allowframebreaks]{Tài liệu tham khảo}
		\begin{thebibliography}{99}
			
			\bibitem{kruger2017}
			T. Krüger, H. Kusumaatmaja, A. Kuzmin, O. Shardt, G. Silva, and E. M. Viggen,
			\emph{The Lattice Boltzmann Method: Principles and Practice},
			Springer, 2017.
			
			\bibitem{chen1998}
			S. Chen and G. D. Doolen,
			``Lattice Boltzmann Method for Fluid Flows,''
			\emph{Annual Review of Fluid Mechanics}, vol. 30, pp. 329–364, 1998.
			
			\bibitem{succi2001}
			S. Succi,
			\emph{The Lattice Boltzmann Equation for Fluid Dynamics and Beyond},
			Oxford University Press, 2001.
			
			\bibitem{palabos}
			FlowKit Ltd.,
			``Palabos: Parallel Lattice Boltzmann Solver,''
			\url{http://www.palabos.org/}, 2024.
			
			\bibitem{lbmethod}
			``The Lattice Boltzmann Method,''
			\url{http://www.lbmethod.org/}, truy cập 2026.
			
			\bibitem{mpi}
			Message Passing Interface Forum,
			``MPI: A Message-Passing Interface Standard,''
			\url{https://www.mpi-forum.org/}, Version 4.0, 2021.
			
		\end{thebibliography}
	\end{frame}
	
	% =================================================================
	% LIÊN HỆ
	% =================================================================
	\hustcontactpage{\hfill THÔNG TIN LIÊN HỆ \hfill }{
		\textbf{Nhóm thực hiện:} Nhóm 1 \\
		\textbf{Môn học:} IT5408 - Tính toán Hiệu năng Cao \\
		\textbf{Giảng viên hướng dẫn:} TS. Nguyễn Quốc Tuấn \\
		
		\vspace{1cm}
		
		\textbf{Source code:} 
		\begin{itemize}
			\item Serial: \texttt{lbm\_serial.c}
			\item MPI: \texttt{lbm\_mpi.c}
			\item Benchmark: \texttt{benchmark.sh}
			\item Visualization: \texttt{visualize\_results.py}
		\end{itemize}
		
		\vspace{1cm}
		\textit{Xin chân thành cảm ơn Thầy và các bạn đã lắng nghe!}
	}
	
	% Trang kết thúc
	\hustthankyou
	
\end{document}